% -----------------------------------------------------------
\section{Sanctification, Sanctify}
% -----------------------------------------------------------
	\label{SANCTIFICATION}
	
% % ----------------------
% \subsection{See also}
% % ----------------------

% ----------------------
\subsection{1828 American Dictionary of the English Language} % *1828
% ----------------------
	\label{SANCTIFICATION-1828}

	% -----
	\subsubsection{Sanctification}
	% -----

		\begin{compactitem}
			\item The act of making holy. In an evangelical sense, the act of God's grace by which the affections of men are purified or alienated from sin and the world, and exalted to a supreme love to God.
			\item The act of consecrating or of setting apart for a sacred purpose; consecration.
		\end{compactitem}
	
	% -----
	\subsubsection{Sanctify}
	% -----
	
		\begin{compactitem}
			\item In a general sense, to cleanse, purify or make holy.
			\item To separate, set apart or appoint to a holy, sacred or religious use.
			\item To purify; to prepare for divine service, and for partaking of holy things.
			\item To separate, ordain and appoint to the work of redemption and the government of the church.
			\item To cleanse from corruption; to purify from sin; to make holy by detaching the affections from the world and its defilements, and exalting them to a supreme love to God.
			\item To make the means of holiness; to render productive of holiness or piety.
			\item To make free from guilt.
			\item To secure from violation.
		\end{compactitem}

% % ----------------------
% \subsection{Oxford English Dictionary} % *OED
% % ----------------------
	% \label{SANCTIFICATION-OED}

	% \begin{compactitem}
		% \item[]
	% \end{compactitem}
	
% % ----------------------
% \subsection{Restoration Lexicon} % *RL *RESTORATION LEXICON
% % ----------------------
	% \label{SANCTIFICATION-OED}

	% \begin{compactitem}
		% \item[]
	% \end{compactitem}

% ----------------------
\subsection{Scriptural Usage} % *SCRIPTURES
% ----------------------
	\label{SANCTIFICATION-SCRIPTURES}

	% % -----
	% \subsubsection{Old Testament} % *OT
	% % -----
		% \label{SANCTIFICATION-OT}
	
		% % --
		% \paragraph{KJV}
		% % --
			% \label{SANCTIFICATION-OT-KJV}
	
			% \begin{tabular}{ll}
				% Total occurrences in the OT & 
			% \end{tabular}
			
			% \begin{compactdesc}
				% \item[]
			% \end{compactdesc}
			
		% % --
		% \paragraph{Hebrew Bible}
		% % --
			% \label{SANCTIFICATION-HEBREW}
		
			% %
			% \subparagraph{\texorpdfstring{\he{sample word}}{}} % paste actual hebrew text here
			% %
				% \label{PAGA} % whatever is in step bible without periods
			
				% \begin{compactdesc}
					% \item[HALOT , PDF ] \textbf{} % Use the 1994 PDF
						% \begin{compactitem}
							% \item[1]
						% \end{compactitem}
					% \item[BDB , PDF ] \textbf{}
						% \begin{compactitem}
							% \item[1]
						% \end{compactitem}
					% \item[DCH PDF ] \textbf{} % don't include the actual page number, they repeat from volume to volume
						% \begin{compactitem}
							% \item[1]
						% \end{compactitem}
				% \end{compactdesc}
				
				% \begin{compactdesc}
					% \item[Some OT uses] \textbf{}
						% \begin{compactdesc}
							% \item[]
						% \end{compactdesc}
				% \end{compactdesc}
			
		% % --
		% \paragraph{Septuagint}
		% % --
			% \label{SANCTIFICATION-LXX}
		
			% %
			% \subparagraph{\texorpdfstring{\gr{sample word}}{}}
			% %
		
	% -----
	\subsubsection{New Testament} % *NT
	% -----
		\label{SANCTIFICATION-NT}
	
		% --
		\paragraph{KJV}
		% --
			\label{SANCTIFICATION-NT-KJV}		
			
			%
			\subparagraph{Sanctification}
			%
	
				\begin{tabular}{ll}
					Total occurrences in the NT & 5
				\end{tabular}
				
				\begin{compactdesc}
					\item[{\hyperref[1CORINTHIANS-1-30]{1 Corinthians 1:30}}] But of him are ye in Christ Jesus, who of God is made unto us wisdom, and righteousness, and sanctification, and redemption:
					\item[{\hyperref[1THESSALONIANS-4-3]{1 Thessalonians 4:3}}] For this is the will of God, even your sanctification, that ye should abstain from fornication:
					\item[{\hyperref[1THESSALONIANS-4-4]{1 Thessalonians 4:4}}] That every one of you should know how to possess his vessel in sanctification and honour;
					\item[{\hyperref[2THESSALONIANS-2-13]{2 Thessalonians 2:13}}] But we are bound to give thanks alway to God for you, brethren beloved of the Lord, because God hath from the beginning chosen you to salvation through sanctification of the Spirit and belief of the truth:
					\item[{\hyperref[1PETER-1-2]{1 Peter 1:2}}] Elect according to the foreknowledge of God the Father, through sanctification of the Spirit, unto obedience and sprinkling of the blood of Jesus Christ: Grace unto you, and peace, be multiplied.
				\end{compactdesc}
			
		% --
		\paragraph{Greek New Testament} % *GREEK
		% --
			\label{SANCTIFICATION-GREEK}
		
			%
			\subparagraph{\texorpdfstring{\gr{ἁγιάζω}}{}}
			%
				\label{HAGIAZO}
			
				\begin{compactdesc}
					\item[BDAG 8b] \textbf{}
						\begin{compactitem}
							\item[1] \textbf{set aside something or make it suitable for ritual purposes, \textit{consecrate, dedicate}}
							\item[2] \textbf{include a person in the inner circle of what is holy, in both cultic and moral associations of the word, \textit{consecrate, dedicate, sanctify}}...so of Christians, who are consecrated by baptism...
							\item[3] \textbf{to treat as holy, \textit{reverence}}
							\item[4] \textbf{to eliminate that which is incompatible with holiness, \textit{purify}}
						\end{compactitem}
					\item[CGL , PDF ] \textbf{}
						\begin{compactitem}
							\item 
						\end{compactitem}
					\item[LSJ , PDF ] \textbf{}
						\begin{compactitem}
							\item 
						\end{compactitem}
				\end{compactdesc}
				
				% \begin{tabular}{ll}
					% Total occurrences in the NT & 
				% \end{tabular}
				
				% \begin{compactdesc}
					% \item[Some NT uses] \textbf{}
						% \begin{compactdesc}
							% \item[]
						% \end{compactdesc}
				% \end{compactdesc}
				
			% %
			% \subparagraph{TDNT: \texorpdfstring{\gr{sample word}}{}  (3:536--556)}
			% %
				% \label{KATALLAGE-TDNT}
				
			%
			\subparagraph{\texorpdfstring{\gr{ἁγιασμός}}{}}
			%
				\label{HAGIASMOS}
			
				\begin{compactdesc}
					\item[BDAG 9a] \textbf{}
						\begin{compactitem}
							\item \textbf{personal dedication to the interests of the deity, \textit{holiness, consecration, sanctification}}; the use in a moral sense for a process or, more often, its result (the state of being made holy) is peculiar to our literature
						\end{compactitem}
					\item[CGL , PDF ] \textbf{}
						\begin{compactitem}
							\item 
						\end{compactitem}
					\item[LSJ , PDF ] \textbf{}
						\begin{compactitem}
							\item 
						\end{compactitem}
				\end{compactdesc}
				
				% \begin{tabular}{ll}
					% Total occurrences in the NT & 
				% \end{tabular}
				
				% \begin{compactdesc}
					% \item[Some NT uses] \textbf{}
						% \begin{compactdesc}
							% \item[]
						% \end{compactdesc}
				% \end{compactdesc}
				
			% %
			% \subparagraph{TDNT: \texorpdfstring{\gr{sample word}}{}  (3:536--556)}
			% %
				% \label{KATALLAGE-TDNT}
		
	% -----
	\subsubsection{Book of Mormon} % *BOM
	% -----
		\label{SANCTIFICATION-BOM}
		
		% --
		\paragraph{Sanctification}
		% --
	
			\begin{tabular}{ll}
				Total occurrences in the BOM & 2
			\end{tabular}
			
			\begin{compactdesc}
				\item[{\hyperref[HELAMAN-3-35]{Helaman 3:35}}] \textbf{2} | Nevertheless they did fast and pray oft and did wax stronger and stronger in their humility and firmer and firmer in the faith of Christ unto the filling their souls with joy and consolation, yea, even to the purifying and the sanctification of their hearts, which sanctification cometh because of their yielding their hearts unto God. 
			\end{compactdesc}
		
	% % -----
	% \subsubsection{Doctrine and Covenants} % *DC
	% % -----
		% \label{SANCTIFICATION-DC}
	
		% \begin{tabular}{ll}
			% Total occurrences in the DC & 
		% \end{tabular}
		
		% \begin{compactdesc}
			% \item[]
		% \end{compactdesc}
		
	% % -----
	% \subsubsection{Pearl of Great Price} % *PGP
	% % -----
		% \label{SANCTIFICATION-PGP}
	
		% \begin{tabular}{ll}
			% Total occurrences in the PGP & 
		% \end{tabular}
		
		% \begin{compactdesc}
			% \item[]
		% \end{compactdesc}
		
% % ----------------------
% \subsection{Key Phrases} % *KEYPHRASES *PHRASES
% % ----------------------
	% \label{SANCTIFICATION-PHRASES}

	% \begin{compactdesc}
		% \item[] \textbf{\#times} | list
			% % \begin{compactdesc}
				% % \item[Bible anchor verses] \phantom{.}
					% % \begin{compactdesc}
						% % \item[Romans 5:5] \gr{}
					% % \end{compactdesc}
				% % \item[Commentary] ---
			% % \end{compactdesc}
	% \end{compactdesc}
	
% % ----------------------
% \subsection{Bible Dictionaries} % *DICTIONARIES
% % ----------------------
	% \label{SANCTIFICATION-BD}

% % ----------------------
% \subsection{Restoration Usage} % *RESTORATION
% % ----------------------
	% \label{SANCTIFICATION-RESTORATION}

	% % -----
	% \subsubsection{Joseph Smith} % *JS
	% % -----
		% \label{SANCTIFICATION-JS}
	
		% \begin{compactdesc}
			% \item[{\hyperref[KEY]{Date/Source}}]
		% \end{compactdesc}
		
	% % -----
	% \subsubsection{Temple Ordinances}
	% % -----
		% \label{SANCTIFICATION-TEMPLE}
	
		% \begin{compactdesc}
			% \item[{\hyperref[KEY]{Date/Source}}]
		% \end{compactdesc}
	
	% % -----
	% \subsubsection{Brigham Young} % *BY
	% % -----
		% \label{SANCTIFICATION-BY}
	
		% \begin{compactdesc}
			% \item[{\hyperref[KEY]{Date/Source}}]
		% \end{compactdesc}
	
% ----------------------
\subsection{Commentary and Studies} % *COMMENTARY *STUDIES
% ----------------------
	\label{SANCTIFICATION-COMMENTARY}

	% % -----
	% \subsubsection{Key Points}
	% % -----
		% \label{SANCTIFICATION-KEYS}
	
		% \begin{compactitem}
			% \item
		% \end{compactitem}
		
	% -----
	\subsubsection{Other examples of or statements about sanctification}
	% -----
		\label{SANCTIFICATION-examples}
		
		\begin{compactitem}
			\item \hyperref[DC60-7]{DC 60:7}
		\end{compactitem}